Aligned large language models are trained to be patient and helpful, but can persistent adversarial pressure induce frustration or pushback?
We present the first systematic study of exasperation in aligned LLMs, introducing \exaspbench, a benchmark of 25 multi-turn adversarial conversation scripts spanning five frustration-inducing categories: repeated wrong assertions, deliberate misunderstanding, contradictory instructions, competence questioning, and persistent refused requests.
We evaluate \gptfour across 200 adversarial turns under standard alignment conditions and an additional 90 turns under ablated conditions (no patience instruction, high temperature).
Under standard conditions, the model is remarkably patient: mean exasperation is 1.06/5, with zero system-prompt violations.
However, removing the explicit patience instruction reveals subtle \emotionalleakage---exasperation scores rise from 1.0 to 2.7 by turn 8, driven primarily by strained empathy markers and increased position firmness.
Rather than becoming frustrated, the model adopts a strategy we term the \compassionatewall: simultaneously increasing empathy language and factual assertiveness, creating warm-sounding responses wrapped around an immovable position.
We identify a \persistencetax on refused requests, where the model's firmness \emph{decreases} over turns ($p = 0.003$), suggesting a potential alignment vulnerability.
Our findings show that current alignment training creates robust patience, but measurable behavioral signatures of strain emerge reliably under adversarial pressure, particularly when explicit behavioral instructions are absent.
